\section{Problem, baseline and principal contribution}
\label{sec:scope}

\subsection{What is being accelerated}

Unknot recognition asks whether an input knot diagram represents a circle
bounding an embedded disk in three-space. The ProveIt programme has developed
several exact certificate paths and a substantial compressed normal-surface
infrastructure. This continuation concerns the latter: given a validated
triangulated source and a supplied normal surface, recognize the homeomorphism
types of the components after cutting, without expanding the potentially huge
number of parallel normal disks. It addresses a terminal subproblem in a
possible hierarchy algorithm. It does not add a new complete scheduler for knot
diagrams or certify that an arbitrary supplied triangulation is a particular
diagram's exterior.

The distinction is algorithmically significant. A normal vector can have only
$7t$ binary coordinates while representing exponentially many disks. An
operation polynomial in the number of represented disks can therefore be
exponential in input size. The existing native constructors use finite source
charts, interval pairings, binary families and source-authenticated orbit
certificates to avoid that expansion. The present work uses their output to
construct a particularly small \emph{active} group presentation. Large families
of isolated balls remain exact integer multiplicities.

The inspected repository is pinned at commit~\cite{RepoSnapshot}
\begin{center}\small\ttfamily
9bacc258bff4b34b4a92b6fa71d290ba5e60ea26
\end{center}
The newest incoming deferred-forests package is present at this revision but
its four-module overlay has not been absorbed wholesale into the production
API~\cite{DeferredForests}. We preserve that distinction when comparing methods.
The new patch applies to this precise production baseline; the deferred
comparison runs in an isolated runtime with the four attributed incoming files.

\subsection{Why a group terminal is the next useful step}

Report 102 established the exact role of unsigned column-leaf cancellation
and exhibited geometric stopping sets with a linear deferral
requirement~\cite{Report102}. The subsequent incoming work removed the need to
enumerate primal forests in the fixed raw positive model and increased
whole-cut coverage from 102 to 111 cases on a retained 120-case corpus. It also
identified a remaining handlebody cut among the refusals and suggested a
free-group terminal under an irreducibility hypothesis~\cite{DeferredForests}.

This article makes that suggestion operational and strengthens its logical
contract. The maintained source validator does not prove irreducibility. We
instead combine exact freeness with the independently authenticated source
Euler characteristic. This detects precisely the second homology discarded
by algebraic simplification. The resulting theorem avoids both an
irreducibility assumption and a separate boundary-component matching step.
The finite compiler reads all source two-handles and all disk-bundle annuli;
it is a presentation of the \emph{whole cut}, not merely the leftover core.

The new move class uses signs and cyclic words. A defining generator can occur
in many other relators, provided it occurs once in its donor; solving for it
and substituting is an exact Tietze transformation. Such substitutions can
produce cancellations invisible to unsigned incidence. This is useful because
the preceding forest theorem already makes repeated primal-tree exploration
unproductive inside its stated raw model. Changing the algebraic move class
addresses a distinct obstruction.

\subsection{Status of the target complexity}

Two author-hosted 2021 slide versions state different quantitative forms of
the announced quasi-polynomial bound: $n^{O(\log n)}$ in the February slides
and $2^{O((\log n)^3)}=n^{O((\log n)^2)}$ in the Oxford handout
\cite{Lackenby2021,LackenbyOxford2021}. Both are quasi-polynomial; their
logarithmic exponents should not be silently identified.

The July 2026 paper supplies a hierarchy framework and compressed cutting
results. Its Section 9 does not estimate the complete runtime; Proposition
10.2 gives a hierarchy length bounded by $4c_q(\mathcal H)$, and Theorem 14.4
gives polynomial compressed-cut processing in handle count and logarithmic
extended weight~\cite{Lackenby2026}. These are valuable ingredients, but do not
by themselves supply the required global search bound. We use the announced
quasi-polynomial regime as a research target, not as a proved property of the
current repository or of this extension.

The older normal-surface complexity results explain why binary compression is
the appropriate starting point. Hass--Lagarias--Pippenger provide classical
complexity bounds and NP certificates for knot problems~\cite{HLP1999}.
Agol--Hass--Thurston give a polynomial interval-orbit algorithm and compressed
normal-surface topology computations~\cite{AgolHassThurston}. Neither result
eliminates the need to control the number of candidate surfaces, the marking
carried through a hierarchy, or the number of expensive choices.

\subsection{Claim ledger}

Table~\ref{tab:claim-ledger} separates the contribution types. The source
compiler and checker are ordinary exact Python programs, not a formalized
proof in a proof assistant. The mathematical proofs are written below; tests
and external comparisons challenge their implementation. No statistical
agreement is substituted for a topological theorem.

\begin{table}[htbp]
\centering\small
\begin{tabular}{P{0.25\linewidth}P{0.32\linewidth}P{0.34\linewidth}}
\toprule
Result & Status & Scope \\
\midrule
Freeness--Euler terminal & Proved from classical topology & Every compact orientable component has nonempty boundary \\
Disk-base presentation & Proved and implemented & Authenticated source; all bases product disks \\
Signed singleton reduction & Proved and implemented & Polynomial, one-sided, incomplete \\
Relator/three-handle ledger & Proved & Exact $b_2$ on completion; lower bound earlier \\
Figure-eight improvement & Source-bound, replayed experiment & One supplied cut, not an unknot verdict \\
Global branch-budget bounds & Conditional theorems & Precise geometric and completeness hypotheses remain open \\
Generic quasi-polynomial recognizer & Open & Neither the current code nor this article establishes it \\
\bottomrule
\end{tabular}
\caption{Logical status of the principal statements.}
\label{tab:claim-ledger}
\end{table}

\subsection{Notation and input contract}

Let $M$ be the source three-manifold and $S$ the supplied normal surface. Write
$W=M\cut S$ for the compact manifold obtained by cutting, rather than for the
set-theoretic complement without its new boundary. The native input is finite,
compact, connected and orientable, with one torus boundary; $S$ satisfies the
normal matching and quadrilateral conditions. Thus $\chi(M)=0$ and
$\chi(W)=\chi(S)$, though the new compiler uses the independently reconstructed
additive source Euler summary. Every component of $W$ has nonempty boundary.

We use $t$ for source tetrahedra, $h_j$ for finite core $j$-handles, $p$ for core
ports, $a,m,L$ for active generators, relators and total literal signed word
length, $c$ for cut components, and $r$ for the sum of their certified free
ranks. All Betti numbers have rational coefficients. A ``free-group
certificate'' proves the entire presented group is free. Computing an
abelianization of the expected rank is insufficient.

The output is \emph{unmarked}: a genus census and, when requested, a proof
packet. It does not supply meridian disks, peripheral transport, a gluing map
to an earlier state, or an essential boundary pattern. In particular a cut of
the figure-eight knot exterior can be a union of handlebodies without the
figure-eight knot being trivial.
